% GENERATED FILE. Edit canonical lessons, then run npm run generate:books.
% canonical-source-hash: fnv1a64:4a6930763cdd7c19
% canonical-lessons: SA-C552-parisilayati, SA-C552-dinacarya, SA-C552-prayasah, SA-C552-prayunkte, SA-C552-upakarakah

\chapter{Trying It Out}
\label{ch:sa-trying-it-out}
\hlchaptermodality{552}

\begin{chapteropening}
\textbf{By the end of this chapter:} \emph{I can say what will happen if I keep trying.}

Complete the last lesson of chapter 552: I can say what will happen if I keep trying.
\end{chapteropening}

\section[\texorpdfstring{\sk{परिशीलयति} (pariśīlayati)}{परिशीलयति (pariśīlayati)}]{\sk{परिशीलयति} (pariśīlayati) — practises regularly}

\label{lesson:SA-C552-parisilayati}



\begin{quote}
Before the new one: say the Sanskrit for detached, then the Sanskrit for lets go of.
\end{quote}

\subsection*{You'll want to know: \sk{परिशीलयति}}
\textbf{\sk{परिशीलयति}} — \emph{pariśīlayati} — \textquotedblleft{}practises regularly\textquotedblright{}.

\begin{quote}
\sk{यदि} \sk{त्वं} \sk{प्रतिदिनं} \sk{परिशीलयसि}, \sk{तर्हि} \sk{त्वं} \sk{कुशलः} \sk{भविष्यसि।}
\end{quote}

\begin{quote}
\emph{yadi tvaṃ pratidinaṃ pariśīlayasi, tarhi tvaṃ kuśalaḥ bhaviṣyasi.}
\end{quote}

\begin{quote}
If you practise every day, you will become good at it.
\end{quote}

Past (he …): \sk{सः} \sk{संगीतं} \sk{परिशीलितवान्}, \emph{saḥ saṅgītaṃ pariśīlitavān}.

\begin{cousinweb}
\textbf{\sk{परिशीलयति}} (\emph{pariśīlayati}) is built on \emph{pari} and the root √śīl, \textquotedblleft{}to practise, to do repeatedly\textquotedblright{}; MW gives \emph{pariśīlayati} as to practise, to use frequently (MW).
\end{cousinweb}

\begin{grammarlens}[title={What you've built}]
One more word for saying what happens if.
\end{grammarlens}

\subsection*{Your turn}
\begin{itemize}
\raggedright
  \item \emph{Say it:} \emph{pariśīlayati}
  \item \emph{Say it:} \emph{pariśīlayati}, once more
  \item \emph{Say it:} \emph{yadi tvaṃ pratidinaṃ pariśīlayasi, tarhi tvaṃ kuśalaḥ bhaviṣyasi.} — If you practise every day, you will become good at it.
  \item \emph{Recall:} say the Sanskrit for detached, then the Sanskrit for lets go of, then say \emph{pariśīlayati} again
\end{itemize}

\begin{tcolorbox}[breakable,colback=teal!4,colframe=teal!35!black,title={Before you move on}]
What does \sk{परिशीलयति} mean? (\textquotedblleft{}Practises regularly\textquotedblright{}.) Say it once more.
\end{tcolorbox}
\section[\texorpdfstring{\sk{दिनचर्या} (dinacaryā)}{दिनचर्या (dinacaryā)}]{\sk{दिनचर्या} (dinacaryā) — a daily routine}

\label{lesson:SA-C552-dinacarya}



\begin{quote}
Before the new one: say the Sanskrit for lets go of, then the Sanskrit for practises regularly.
\end{quote}

\subsection*{You'll want to know: \sk{दिनचर्या}}
\textbf{\sk{दिनचर्या}} — \emph{dinacaryā} — \textquotedblleft{}a daily routine\textquotedblright{}.

\begin{quote}
\sk{यदि} \sk{दिनचर्या} \sk{उत्तमा} \sk{अस्ति}, \sk{तर्हि} \sk{स्वास्थ्यम्} \sk{अपि} \sk{उत्तमं} \sk{भवति।}
\end{quote}

\begin{quote}
\emph{yadi dinacaryā uttamā asti, tarhi svāsthyam api uttamaṃ bhavati.}
\end{quote}

\begin{quote}
If your daily routine is good, your health is good too.
\end{quote}

A feminine noun.

\begin{cousinweb}
\textbf{\sk{दिनचर्या}} (\emph{dinacaryā}) is built from \emph{dina}, \textquotedblleft{}a day\textquotedblright{}, and \emph{caryā}, \textquotedblleft{}going about, conduct\textquotedblright{}, from the root √car, \textquotedblleft{}to move\textquotedblright{}; MW gives \emph{dinacaryā} as daily work (MW).
\end{cousinweb}

\begin{grammarlens}[title={What you've built}]
One more word for saying what happens if.
\end{grammarlens}

\subsection*{Your turn}
\begin{itemize}
\raggedright
  \item \emph{Say it:} \emph{dinacaryā}
  \item \emph{Say it:} \emph{dinacaryā}, once more
  \item \emph{Say it:} \emph{yadi dinacaryā uttamā asti, tarhi svāsthyam api uttamaṃ bhavati.} — If your daily routine is good, your health is good too.
  \item \emph{Recall:} say the Sanskrit for lets go of, then the Sanskrit for practises regularly, then say \emph{dinacaryā} again
\end{itemize}

\begin{tcolorbox}[breakable,colback=teal!4,colframe=teal!35!black,title={Before you move on}]
What does \sk{दिनचर्या} mean? (\textquotedblleft{}A daily routine\textquotedblright{}.) Say it once more.
\end{tcolorbox}
\section[\texorpdfstring{\sk{प्रयासः} (prayāsaḥ)}{प्रयासः (prayāsaḥ)}]{\sk{प्रयासः} (prayāsaḥ) — exertion}

\label{lesson:SA-C552-prayasah}



\begin{quote}
Before the new one: say the Sanskrit for practises regularly, then the Sanskrit for a daily routine.
\end{quote}

\subsection*{You'll want to know: \sk{प्रयासः}}
\textbf{\sk{प्रयासः}} — \emph{prayāsaḥ} — \textquotedblleft{}exertion\textquotedblright{}.

\begin{quote}
\sk{यदि} \sk{प्रयासः} \sk{सत्यः} \sk{अस्ति}, \sk{तर्हि} \sk{फलम्} \sk{अवश्यं} \sk{भवति।}
\end{quote}

\begin{quote}
\emph{yadi prayāsaḥ satyaḥ asti, tarhi phalam avaśyaṃ bhavati.}
\end{quote}

\begin{quote}
If the effort is real, a result surely comes.
\end{quote}

A masculine noun.

\begin{cousinweb}
\textbf{\sk{प्रयासः}} (\emph{prayāsaḥ}) is built on \emph{pra} and the root √yas, \textquotedblleft{}to heat oneself, to exert oneself\textquotedblright{}; MW gives \emph{prayāsa} as exertion, effort, pains, trouble (MW).
\end{cousinweb}

\begin{grammarlens}[title={What you've built}]
One more word for saying what happens if.
\end{grammarlens}

\subsection*{Your turn}
\begin{itemize}
\raggedright
  \item \emph{Say it:} \emph{prayāsaḥ}
  \item \emph{Say it:} \emph{prayāsaḥ}, once more
  \item \emph{Say it:} \emph{yadi prayāsaḥ satyaḥ asti, tarhi phalam avaśyaṃ bhavati.} — If the effort is real, a result surely comes.
  \item \emph{Recall:} say the Sanskrit for practises regularly, then the Sanskrit for a daily routine, then say \emph{prayāsaḥ} again
\end{itemize}

\begin{tcolorbox}[breakable,colback=teal!4,colframe=teal!35!black,title={Before you move on}]
What does \sk{प्रयासः} mean? (\textquotedblleft{}Exertion\textquotedblright{}.) Say it once more.
\end{tcolorbox}
\section[\texorpdfstring{\sk{प्रयुंक्ते} (prayuṅkte)}{प्रयुंक्ते (prayuṅkte)}]{\sk{प्रयुंक्ते} (prayuṅkte) — makes use of}

\label{lesson:SA-C552-prayunkte}



\begin{quote}
Before the new one: say the Sanskrit for a daily routine, then the Sanskrit for exertion.
\end{quote}

\subsection*{You'll want to know: \sk{प्रयुंक्ते}}
\textbf{\sk{प्रयुंक्ते}} — \emph{prayuṅkte} — \textquotedblleft{}makes use of\textquotedblright{}.

\begin{quote}
\sk{यदि} \sk{भवान्} \sk{संगणकं} \sk{प्रयुंक्ते}, \sk{तर्हि} \sk{कार्यं} \sk{शीघ्रं} \sk{भविष्यति।}
\end{quote}

\begin{quote}
\emph{yadi bhavān saṅgaṇakaṃ prayuṅkte, tarhi kāryaṃ śīghraṃ bhaviṣyati.}
\end{quote}

\begin{quote}
If you use a computer, the work will be quick.
\end{quote}

Past (he …): \sk{सः} \sk{संगणकं} \sk{प्रयुक्तवान्}, \emph{saḥ saṅgaṇakaṃ prayuktavān}.

\begin{cousinweb}
\textbf{\sk{प्रयुंक्ते}} (\emph{prayuṅkte}) is built on \emph{pra}, \textquotedblleft{}forth\textquotedblright{}, and the root √yuj, \textquotedblleft{}to yoke\textquotedblright{}; MW gives \emph{prayuj} as to use, to employ, to practise (MW).
\end{cousinweb}

\begin{grammarlens}[title={What you've built}]
One more word for saying what happens if.
\end{grammarlens}

\subsection*{Your turn}
\begin{itemize}
\raggedright
  \item \emph{Say it:} \emph{prayuṅkte}
  \item \emph{Say it:} \emph{prayuṅkte}, once more
  \item \emph{Say it:} \emph{yadi bhavān saṅgaṇakaṃ prayuṅkte, tarhi kāryaṃ śīghraṃ bhaviṣyati.} — If you use a computer, the work will be quick.
  \item \emph{Recall:} say the Sanskrit for a daily routine, then the Sanskrit for exertion, then say \emph{prayuṅkte} again
\end{itemize}

\begin{tcolorbox}[breakable,colback=teal!4,colframe=teal!35!black,title={Before you move on}]
What does \sk{प्रयुंक्ते} mean? (\textquotedblleft{}Makes use of\textquotedblright{}.) Say it once more.
\end{tcolorbox}
\section[\texorpdfstring{\sk{उपकारकः} (upakārakaḥ)}{उपकारकः (upakārakaḥ)}]{\sk{उपकारकः} (upakārakaḥ) — helpful}

\label{lesson:SA-C552-upakarakah}



\begin{quote}
Before the new one: say the Sanskrit for exertion, then the Sanskrit for makes use of.
\end{quote}

\subsection*{You'll want to know: \sk{उपकारकः}}
\textbf{\sk{उपकारकः}} — \emph{upakārakaḥ} — \textquotedblleft{}helpful\textquotedblright{}.

\begin{quote}
\sk{यदि} \sk{उपायः} \sk{उपकारकः} \sk{अस्ति}, \sk{तर्हि} \sk{अहं} \sk{तं} \sk{करिष्यामि।}
\end{quote}

\begin{quote}
\emph{yadi upāyaḥ upakārakaḥ asti, tarhi ahaṃ taṃ kariṣyāmi.}
\end{quote}

\begin{quote}
If the method is helpful, I will use it.
\end{quote}

\begin{cousinweb}
\textbf{\sk{उपकारकः}} (\emph{upakārakaḥ}) is built on \emph{upa}, \textquotedblleft{}towards\textquotedblright{}, and the root √kṛ, \textquotedblleft{}to do\textquotedblright{}, with the agent ending \emph{-aka}, like \emph{upakāra}, \textquotedblleft{}a favour\textquotedblright{}; MW gives \emph{upakāraka} as doing a service, helping (MW).
\end{cousinweb}

\begin{grammarlens}[title={What you've built}]
One more word for saying what happens if.
\end{grammarlens}

\subsection*{Your turn}
\begin{itemize}
\raggedright
  \item \emph{Say it:} \emph{upakārakaḥ}
  \item \emph{Say it:} \emph{upakārakaḥ}, once more
  \item \emph{Say it:} \emph{yadi upāyaḥ upakārakaḥ asti, tarhi ahaṃ taṃ kariṣyāmi.} — If the method is helpful, I will use it.
  \item \emph{Recall:} say the Sanskrit for exertion, then the Sanskrit for makes use of, then say \emph{upakārakaḥ} again
\end{itemize}

\begin{tcolorbox}[breakable,colback=teal!4,colframe=teal!35!black,title={Before you move on}]
What does \sk{उपकारकः} mean? (\textquotedblleft{}Helpful\textquotedblright{}.) Say it once more.
\end{tcolorbox}
